% !TEX program = xelatex
% ============================================================
% 玉器门店顾客服务与经营管理系统 —— 需求调研报告
% 课程：系统分析与设计    编译：xelatex
% ============================================================
\documentclass[12pt, a4paper]{ctexart}

\usepackage[top=2.4cm, bottom=2.4cm, left=2.6cm, right=2.6cm]{geometry}
\usepackage{graphicx}
\usepackage{booktabs}
\usepackage{longtable}
\usepackage{array}
\usepackage{tabularx}
\usepackage{enumitem}
\usepackage{xcolor}
\usepackage{caption}
\usepackage{float}
\usepackage{hyperref}

\hypersetup{colorlinks=true, linkcolor=black, urlcolor=blue, citecolor=black}
\captionsetup{font=small, labelfont=bf, labelsep=quad}
\counterwithin{table}{section}
\counterwithin{figure}{section}
\renewcommand{\thetable}{\thesection-\arabic{table}}

\title{玉器门店顾客服务与经营管理系统\\[6pt] \large 需求调研报告（调研对象：新疆玉器门店老板）}
\author{姓名：\underline{张昊轩}\quad 学号：\underline{8102240429}\quad 班级：\underline{信管 T2401}}
\date{课程：系统分析与设计 \quad 指导教师：毕文杰 \quad 2026 年 10 月}

\begin{document}

\maketitle
\thispagestyle{empty}

\section{调研概况}

\begin{table}[H]
\centering
\caption{调研概况}
\begin{tabularx}{\textwidth}{@{}p{2.8cm} X@{}}
\toprule
项目 & 内容 \\
\midrule
调研时间 & 2026 年 10 月 1 日 \\
调研方式 & 对门店老板当面访谈；辅以公众号「长生玉鉴」已发文章三篇（《和田玉科普：聊聊被神话的丁达尔效应》《马鬃山寻玉。》《实战：新一代高仿戈壁料，逛摊避坑心得》）作为背景材料 \\
调研对象 & 新疆的玉器门店老板 \\
调研目的 & 为系统的愿景定义、业务建模与需求分析提供事实依据 \\
\bottomrule
\end{tabularx}
\end{table}

\section{门店规模与经营现状}

\begin{itemize}[leftmargin=1.6em, itemsep=2pt]
  \item 门店在新疆，在柜 4～5 人；在柜货品约 3000～4000 件，老板说“这个数并不准，很多是小件，有些货还能拆”。
  \item 在柜货品总价约 2000～3000 万元，其中相当一部分是同行调来的货，不都是门店自有。
  \item 成交单数不固定；客单价跨度极大，低的一两千，高的到六位数。
  \item 毛利一年一百多万（老板口述时先说“净利润”，后更正为毛利）；门店有合伙人、要发工资、要养几个人，“提销售额还是很重要的”。
  \item 老板自述：“成本这一块主要靠眼力，所以很难通过系统去提销售额，这个我很清楚。”
\end{itemize}

\section{记账现状与最大痛点}

\begin{itemize}[leftmargin=1.6em, itemsep=2pt]
  \item 现状：纸质台账 + Excel；调货全靠微信——“在图片上面写这个价格，然后发给别人，别人拿走”。
  \item 微信记录不能删，但“图片有时候它自动就删了，删了之后就会扯皮”。
  \item 最头疼的是调货：“总会丢很多货……记录丢失了就不好再去找当时调出去的货”。调货每天都有很多笔。
  \item 识别同一件货目前靠三样：克重、大体照片、证书。
\end{itemize}

\section{周转周期（愿景指标的现实依据）}

\begin{itemize}[leftmargin=1.6em, itemsep=2pt]
  \item 老板原话：“一件货如果是六位数的，现货可能压个三四年都是有可能的；小几百到一两千的货，可能压一周多一点。”
  \item 结论：高价货压得久、低价货走得快，必须按货品金额加权，否则平均值会掩盖真实积压。
  \item 压周转“我认为是非常需要压的”；销售额提升 10\%、周转周期下降 20\%“就是估的”。
\end{itemize}

\section{货品来源与收料}

\begin{itemize}[leftmargin=1.6em, itemsep=2pt]
  \item 四类来源占比（老板口径）：成品进货约 10\%、回收回流约 20\%、同行调货约 40\%、买料加工约 30\%。
  \item 以买料加工为主、同行调货为辅，成品与回收较少。主做新疆料，对新疆料要求高；切山料和籽料，戈壁料、籽料也卖一些；原石山流水料买得少——“风险比较大”。
  \item 收料两条路：自己在新疆收（逛摊、跑矿区）；固定供货——“会有一些少数民族的兄弟把货送到门店，让我们自己挑料子”。
  \item 通货买法：一次性把对方的货全买下，里面有“货头货尾”（品质高的和低的）。
  \item 供货痛点：“提前是看不到货的，他只会跟我约个大概的时间，可能他来的时候我们也不在，或者他根本没有提前约时间——这都非常常见。”
  \item 真假与产地：“产地尤其靠眼力，现在证书是做不出来产地的”；玉的真假“主要还是要依靠鉴定证书”。老板认为上了系统后“我们的知信度就会很高，相对其他的门店来说”。
\end{itemize}

\section{成交、价格与客户}

\begin{itemize}[leftmargin=1.6em, itemsep=2pt]
  \item 定价：成本价上加 30\%～40\% 毛利定出底价（最低售价）；标价留溢价空间；低于底价要老板点头。
  \item 收款：以全款为主，定金较少——定金是“分期留货，钱付完了再把货拿走”。
  \item 客户结构：老客户多，散客多是老客户介绍来的；没有网络关联的散客很少直接进店。
  \item 客户资料现在基本记在微信里；货品对外展示靠朋友圈和微商相册。
  \item 引流：公众号效果一般，主要靠小红书和抖音，“有人看了小红书专门来店里”。
  \item 到店协调的麻烦：线上客户找老板预约，但“店里销售今天有事情，或者已经约了别的客户，就没办法接待”，老板要两头协调。
  \item 老板态度明确：“看货、预约、预定，那肯定是要跑起来的”——客户在系统上约好，来店直接找对应销售。
\end{itemize}

\section{退货、鉴赏期与调货结算}

\begin{itemize}[leftmargin=1.6em, itemsep=2pt]
  \item 退货：成交后 7 天内、限真假与质量问题，写进收据（小票），系统里也要写明。
  \item 主观不满意：老板说“一般还是要协商一下的，看看能不能换货；实在不行再与销售协商”——退款不办，换货可以谈，而且“和销售的业绩是相关的”。
  \item 鉴赏期：和田玉行规，客户拿到货可以鉴赏 3 天或 7 天再决定；鉴赏期一过就要结账（调货结算的时点）。
  \item 调货结算：调出去的货“当然有记错有扯皮的”，根源就是微信记录会丢。
\end{itemize}

\section{系统形态与边界}

\begin{itemize}[leftmargin=1.6em, itemsep=2pt]
  \item 柜台：以电脑为主，手机也有，店里有网络。
  \item 外部角色（供应商、加工师傅、同行）必须手机访问——老板原话：“他们用电脑肯定不现实的，而且文化程度不高……做成小程序还是更好的。”
  \item 供应商：预约上门时间，上传料子图片。
  \item 加工师傅：老板先电话交代，师傅在系统上提交——做之前什么样、中间上传一次、做完上传一次，然后通知老板去拿或送过来；老板在系统记账、打钱，师傅确认收到。
  \item 调货（老板认为最重要）：“同行能看到我有哪些货可以调给他，他也有哪些货可以调给我，都可以写到系统里面”——先老板与同行之间用，以后可以同行与同行。
  \item 为什么不用现成进销存：“玉器业务不是标准化的，我做一个手串撑死了做 100 条，这 100 条也有不同的定价，长得也很不一样，克重、质量都不同，用现成进销存就要设 100 个不同的项目，非常麻烦。”
  \item 系统先解决三件事：入库出库记账、同行调货、责任划分——“入库出库是老板和销售同时管理的，会记错，但找不着责任人”。
  \item 明确不管：支付（走公账、线下微信支付宝）、物流（顺丰）、鉴定机构（会员制，人去就行，交给销售）。送检由老板在系统提交货品清单，销售每天去一两次把证书做出来；老板不在店时由销售代理收货，“通过系统信息化，最后也有一个责任的划分，会相对比较清楚”。
\end{itemize}

\section{既有系统的能力（客户管理工具）}

老板手上已有一个自研的桌面小工具（Python + SQLite 单机版，含客户、订单、货品、利润、统计五个页面），是需求的直接来源，其产品文档中已固化的做法如下。

\begin{itemize}[leftmargin=1.6em, itemsep=2pt]
  \item 客户档案字段：编号、昵称/姓名、手机号、微信号、来源渠道（下拉枚举：闲鱼 / 微信朋友圈 / 朋友介绍 / 线下 / 直播 / 老客户复购 / 其他）、标签（大客户 / 懂行 / 爽快 / 要还价）、备注。
  \item 检索：按编号、昵称、手机号、微信号、备注、标签任一条件模糊搜索；列表页手机号可一键复制，便于直接加微信。
  \item 客户详情：累计消费、单数、占比、最近成交时间、成交历史、跟进记录。
  \item 核心客户看板：一眼看出贡献营业额大头的核心客户及其历史成交。
  \item 来源渠道分析：按渠道统计成交，判断哪个渠道真正有效。
  \item 货品一物一档：多图挂载、料种、产地、器型、克重、尺寸、证书号，以及“在库 / 已售 / 预留 / 已退货 / 代销中”的状态语义。
  \item 明确不做（老板已决策）：库存预警与补货提醒、客户生日提醒、会员等级与积分优惠券、营销自动化与短信群发、多用户权限与云端同步、Excel 导入。
\end{itemize}

结论：客户管理这块，老板要的不是“会员营销系统”，而是联系方式结构化的客户档案、全字段找得到人、看得出谁是核心客户与哪个渠道有效。

\section{调研结论}

调研结论按主题归纳如下。

\begin{table}[H]
\centering
\small
\caption{调研结论归纳}
\begin{tabularx}{\textwidth}{@{}p{2.6cm} X@{}}
\toprule
主题 & 主要结论 \\
\midrule
经营现状 & 门店在新疆，4～5 人在柜，在柜货品 3000～4000 件，货值约 2000～3000 万元（含大量同行调货）；客单价一两千到六位数；毛利一年一百多万。收料定价靠眼力，系统难以直接提升销售额，主要支撑老客户复购与线上引流到店的转化 \\
记账与痛点 & 纸质台账加 Excel 记账；调货靠微信发图写价，记录会丢、对账扯皮，是当前最大痛点；识别同一件货靠克重、照片与证书 \\
周转周期 & 六位数的货压三四年、几百到一两千的货一周多，周转周期必须按货品金额加权；销售额提升 10\%、周转周期下降 20\% 为拟议目标 \\
货品来源 & 成品进货约 10\%、回收回流约 20\%、同行调货约 40\%、买料加工约 30\%；主做新疆料，以买料加工为主、同行调货为辅；供货方上门时间不定，需要预约与代收机制 \\
成交与价格 & 成本加 30\%～40\% 毛利定出底价，低于底价须老板授权；全款为主、定金分期留货；一物一码以条形码标签落地，成交登记可扫码识别货品 \\
客户与引流 & 客户资料散落在微信；老客户为主、线上引流到店；客户档案需结构化（联系方式、来源渠道、标签、跟进记录），统计核心客户贡献与渠道效果 \\
退货与鉴赏期 & 成交后 7 天内、限真假与质量问题可退，政策写入收据；主观不满意不退款、协商换货；同行调货以客户鉴赏期（3 天或 7 天）届满为结账时点 \\
系统形态 & 柜台电脑端，外部协同用微信小程序，操作步骤尽量少；货品标签用条形码，鉴定证书二维码由机构印制、扫码可查线上证书 \\
系统边界 & 不接入在线支付、不做物流、不对接鉴定机构系统；送检由老板提交货品清单、销售集中办理，老板不在时销售代理收货并留经手人 \\
\bottomrule
\end{tabularx}
\end{table}

\newpage
\appendix
\section{访谈原始记录（老板口述，未作改动）}

\noindent
{\small 说明：以下为被调研人当场口述的逐字记录，保留口语与行业说法，仅对明显的口误加注，供核对事实使用。}

\medskip
\noindent
店里现在一共有 4～5 个人在柜，大概有 3000～4000 件货吧。然后啊，因为有很多是比较小的，还有一些货可以进行拆分，所以货的这个并不准，然后总价大概在 2～3000 万，但是并不都是自己的，其中有很多是进行调货过来的，一个月成交大概有。嗯，这个不好说，因为成交单数是不固定的，客单价也有低有高啊，最低会在一两千，最高的话就是在 6 位数左右。

\medskip
\noindent
现在账就是纸质台账，纸质台账加 excel，然后微信聊天记录是不能删的。因为现在的调货都是通过微信聊天记录去调的，在图片上面写这个价格，然后发给别人，别人拿走，然后就是这样子。啊，现在最让头疼的事情就是调货。他啊总会丢很多货，因为通过微信聊天的话，它还是不准啊，然后会经常丢失一些记录啊。记录丢失了就不好再去找当时调出去的货的这个样子了。我们调货的时候，确认同一件货的方法一般就是克重，克重和大体的一个照片。啊，然后就是证书，证书也很重要。

\medskip
\noindent
然后提销售额，压周转周期，这两个数其实心里也不知道，现在一个月。能卖个就是净利润啊，净利润的话一年可能有个 100 多万吧〔注：此处老板口述为“净利润”，后更正为毛利〕，然后呢。然后这个因为合伙人也比较多，然后呃，再加上发工资，然后档口这几个人的生活，所以。嗯，提销售额还是很重要的。成本这一块，还是主要靠眼力，所以很难通过系统去提销售额。这个我很清楚，压周转周期的话，我认为啊。这个出产周期是非常需要压的，现在一件货价格越高的一般压得越久，所以我肯定是价格越高的想压的越少啊。这就是估的估的这个一个百分比。我一件货如果是 6 位数的，我这个现货可能压个三四年都是有可能的。要是这个小几百到一两千的这种货啊，可能就是压一周左右吧，一周多一点。这样还是没什么问题，主要是大价格这个，所以要加权，用价格去加权。

\medskip
\noindent
然后货的来源就是这四种。啊，成品然后回收的这个回流货，然后同行调货来，然后就是我们买原料自己加工，因为我们主做新疆料，对于新疆料的这个要求是比较高的。成品进货量很少，大概占个 10\% 左右吧。回收回流这个价格便宜，可能占个 20\%。然后呢，同行调货和买料加工，同行调货可能占 40\%。然后剩下百分之三十百百分之 30，它就是一个买料加工了。我们买料加工呢，就是切山料和籽料啊，然后戈壁料和籽料也会卖一些。原石山流水料买的很少，因为它的风险比较大啊。然后以哪种为主，就是还是以买料加工为主，然后同行调货为辅，其他两个相对比较少。

\medskip
\noindent
收料是自己在新疆收的比较多，固定供货的也有，会有一些啊，这个少数民族。少数民族的兄弟过来送货，会送到门店，把货给我们送到门店，然后让我们自己挑料子。这会儿有可能是通货一起买，通货的意思就是我们呃一次性把他的所有的货都买下来，然后呢，其中可能有货头货尾，就是品质高的品质低的这样子。啊，然后这个也是一个比较大的问题啊，这是供货这边的问题，因为我们提前是看不到的，他只会跟我们约个大概的时间，可能他来的时候我们也不在，或者是他都没有提前来约时间，这都非常常见。

\medskip
\noindent
从摊上收来的料的话，嗯，真假和产地是主要靠自己的眼力，产地尤其靠眼力，因为现在证书是做不出来产地的，我们要做这个事情的时候，我们要如果有个系统，我认为我们的这个就是。呃，知信度就会很高，就是相对其他的门店来说，因为我们系统会更加的呃方便快捷，对于客户来说。然后玉的真假的话，我们主要还是要依靠鉴定证书。

\medskip
\noindent
成交价怎么定呢？成交价一般我们会在成本价，就是有底价，一般来说啊，一般正常不是捡漏的，不是远低于市场价收来的那种东西，我们会在这个呃基本的价格上面定个 30\%～40\% 的毛利，啊，然后呢。这个毛利率去固定它的一个最低价啊，然后卖货的这个标价是比较高的，相对来说比较高的啊是要有有留有一个溢价空间。当然有一个底价和最低售价了，低于底价需要我点头，需要我一个老板点头。然后，呃，全款和定金大概各占多少？全款为最重要的，就是基本上都是全款，定金也有比较少，可能这个人这个分付出来这么多钱，然后他会分期留这个货，然后等他所有的钱付完了之后，他再把这个货品拿走。啊，这个你也需要做一下用力，这个肯定是。啊，需要写到这个用例规约里面的。

\medskip
\noindent
老客户比较多，散客还是呃老客户去介绍的比较就是多的，因为散客他啊就是没有任何的呃网络关联的情况下，他很少来直接找我买，还有就是线上引流，线上引流还是比较多的，线上引流他。呃，我们会有一些可能没有来过店里的的朋友，然后他会来找我啊，老板来预约。但是店里头可能销售他今天有事情，或者是销售今天啊已经约到了呃，他需要去接待的客户，他就会没有办法去预约。我这边就需要两头去呃，这个协调还是比较麻烦的。啊，然后客户资料现在记在基本上都是记在微信里面嘛。然后啊，我们的货品入库就是用朋友圈和微商相册这两个东西。公众号现在引流效果不是很好，主要是靠小红书和抖音嘛。抖音效果也不是很好啊。然后有人专门看了文章来店里了，就是小红书的啊。

\medskip
\noindent
你这个看货预约预定，那肯定是要跑起来的。你这个如果他想来店里，我可以直接给他发系统，然后他在系统上去约，他就可以直接去找了销售啊，然后和销售去进行一个沟通交流啊，然后呢，退货与调货，这个退货。7 天，这个真假质量这个刚说了啊。这个写进小票啊，当然写进收据里啊。然后，这个我们也会在系统里面去明说不喜欢想退，一般我们还是要协商一下的啊，协商为什么不喜欢，然后看看能不能进行换货。如果实在不行的话，再与销售进行协商，因为这个和销售的业绩是相关的。

\medskip
\noindent
然后同行调货一般。就是每天都有很多很多啊。结账就是客户这边过了鉴赏期啊，和田玉是有个鉴赏期的，就是他拿到货之后可以呃这个鉴赏个 3 天或者 7 天这样的，然后来确定他到底喜不喜欢这个东西啊。就是他建档期一过，然后这边就要结账了，调出去的货当然有记错有扯皮的啊，因为这个微信靠纯靠微信微信记录，微信记录的话，这个图片啥的就是有时候可能它自动就删了，删了之后就会扯皮啊。

\medskip
\noindent
柜台是用电脑，手机也有啊，主要还是靠电脑，然后网络的话是有网络的啊，供应商师傅同行。啊是需要访问的供应商这边就是供应商师傅同行，最好还是用手机去访问他们啊用电脑它不太现实。啊，好多都是散的那种档口柜台，他用电脑肯定不现实的，而且文化程度不高，所以做成小程序，或者是做成啊非常好用的那种链接，还是小小程序吧，小程序还是更好的。嗯。

\medskip
\noindent
供应商这边就是他预约一下什么时候来上门，然后他上传一下他料子的图片啊。然后师傅这边就是我这个打电话跟他说，然后他在系统上提交一下我这个哪些图片做之前是什么样的，就切掉切完了之后做之前什么样的，然后做到中间给我上传一次。然后做完了之后，然后给我上传一次，完了之后，他就可以告诉我可以去拿了，我就去拿，或者他给我送过来，他在上面提交。提交完了之后，我给他把钱通过系统去进行一个呃记账。记完账之后，我给他打钱，然后他确定收到就行了。

\medskip
\noindent
然后呢，调货这个更重要，调货是最重要的，这个事情就同行调货，他能看到我有哪些货可以调给他，他也有哪些货可以调给我，他都可以写到系统里面。然后，这个系统刚开始可以就是我和同行之间去使用，后面也可以同行与同行之间去使用。现在就是我和同行使用。

\medskip
\noindent
现成的进销存，它可能是不太相关，我们的这个玉器业务，玉器业务它不是标准化了，我可能一个手串我撑死了，我做了 100 条。这 100 条它也有不同的定价，它长得也很不一样，它的克重，它的质量都不同的。呃，我用现存的进销，现成的进销存的话，我就需要呃设置 100 个不同的项目。啊，这都非常的麻烦。我们现在系统肯定要先解决记账的问题，记账入库出库。现在记账入库，出库是老板和销售同时去管理的，然后呢也会出现记错的情况，但是找不着责任人。然后呢，还有就是同行的调货，调货这个太混乱了，必须要管。

\medskip
\noindent
然后哪些是明确不让管，这支付肯定不用你管啊。支付我们还是走公账，走公账和线下支付，线下的一个微信，支付宝的一个支付。物流的话，我们找顺丰公司经。然后鉴定机构你也甭管鉴定机构，因为它是一个会员制的，然后只要人去就可以了，所以就是交给销售去做啊，其实就是我在系统上找这个销售啊销售这边。呃，看到我提交了这个哪些货品需要去做证书，或他每天去一到 2 次把所有证书做出来啊就可以了。然后出现收货的时候收货这种情况，然后证书也是由销售去做啊，然后这都是啊，通过系统可以去就是信息化的啊，你明白吧，因为很多时候可能收货的时候啊，我把货收回来了，但是我人并不在店里，或者说他人来店里了。我看货的料还不错，但是我人还没回来，就是呃送货的人来店里了，我人还没回来。这种时候也需要销售去进行一个代理，所以呢，我们还是把它信息化会比较好一些。而且通过系统信息化的话，最后也有一个责任的划分会相对比较清楚嗯。

\end{document}
